\documentclass[11pt,a4paper,notitlepage]{report}
% Shared preamble for the raw chapter-4 write-ups under docs/report/insights/.
%
% Each experiment_N/experiment_N.tex is a standalone, compilable LaTeX report
% written BEFORE anything goes into final_report/chapters/chapter_4_*.tex.
% It uses the final report's fonts, units and macros so that passages, tables
% and figures lift into the chapter without reformatting.
%
%   \documentclass[11pt,a4paper]{report}
%   % Shared preamble for the raw chapter-4 write-ups under docs/report/insights/.
%
% Each experiment_N/experiment_N.tex is a standalone, compilable LaTeX report
% written BEFORE anything goes into final_report/chapters/chapter_4_*.tex.
% It uses the final report's fonts, units and macros so that passages, tables
% and figures lift into the chapter without reformatting.
%
%   \documentclass[11pt,a4paper]{report}
%   % Shared preamble for the raw chapter-4 write-ups under docs/report/insights/.
%
% Each experiment_N/experiment_N.tex is a standalone, compilable LaTeX report
% written BEFORE anything goes into final_report/chapters/chapter_4_*.tex.
% It uses the final report's fonts, units and macros so that passages, tables
% and figures lift into the chapter without reformatting.
%
%   \documentclass[11pt,a4paper]{report}
%   \input{../common/preamble}
%   \begin{document} ...
%
% Figures and tables are NOT copied: \graphicspath and \reporttable resolve into
% final_report/figures/ and final_report/tables/, so the raw write-up always
% shows the same artefacts as the chapter.  Compile with:  latexmk -pdf <file>

\usepackage[T1]{fontenc}
\usepackage[utf8]{inputenc}
\usepackage{mathptmx}
\usepackage[scaled=0.92]{helvet}
\usepackage{courier}
\usepackage[a4paper,left=30mm,right=20mm,top=25mm,bottom=25mm]{geometry}
\usepackage{microtype}
\usepackage{booktabs}
\usepackage{tabularx}
\usepackage{longtable}
\usepackage{array}
\usepackage{multirow}
\usepackage{amsmath}
\usepackage{amssymb}
\usepackage{ifthen}
\usepackage{siunitx}
\usepackage{graphicx}
\usepackage{placeins}
\usepackage{caption}
\usepackage{enumitem}
\usepackage[table]{xcolor}
\usepackage{rotating}
\PassOptionsToPackage{hyphens}{url}\usepackage{url}
\usepackage[hidelinks]{hyperref}
\usepackage[nameinlink,noabbrev]{cleveref}

\graphicspath{{../../final_report/}{../../final_report/figures/}}
% \reporttable{name} inputs final_report/tables/name.tex
\newcommand{\reporttable}[1]{\input{../../final_report/tables/#1}}

\newcommand{\repo}[1]{\texttt{\detokenize{#1}}}
\sisetup{detect-all,group-separator={,},range-phrase=--,range-units=single}
\DeclareSIUnit{\pixel}{px}
\DeclareSIUnit{\mebi}{Mi}
\captionsetup{font=small,labelfont=bf}
\setlist{nosep}
\setlength{\parskip}{4pt}
\setlength{\LTcapwidth}{\textwidth}

% ---- result-slot status (RENDERED here, unlike in the final report)
\definecolor{stcompleted}{HTML}{A5D6A7}
\definecolor{stpartial}{HTML}{FFF2CC}
\definecolor{stuntested}{HTML}{FFD8A8}
\definecolor{stblocked}{HTML}{FFCDD2}
\definecolor{stfailed}{HTML}{FFCDD2}
\definecolor{stplanned}{HTML}{E0E0E0}
\newcommand{\slotstatus}[2]{%
  \par\noindent\fcolorbox{black!40}{st#1}{\parbox{\dimexpr\textwidth-2\fboxsep-2\fboxrule}{%
    \textbf{Status: #1.}\ #2}}\par\medskip}
\newcommand{\status}[1]{\textbf{#1}}

% ---- queries: RENDERED in red here so a reader sees every open concern in place
\definecolor{flagred}{HTML}{B00000}
\newcommand{\flag}[1]{\textcolor{flagred}{#1}}
\newcommand{\moved}[1]{#1}
\newenvironment{queried}[1][Questionable]%
  {\par\addvspace{0.6ex}\begingroup\color{flagred}\noindent
   \textbf{[#1]}\ \ignorespaces}%
  {\par\endgroup\addvspace{0.6ex}}

% ---- per-block performance marks (section 4.5 diagrams)
\definecolor{perfgood}{HTML}{A5D6A7}
\definecolor{perfpoor}{HTML}{FFCDD2}
\definecolor{perfnoeffect}{HTML}{D9EAF7}
\definecolor{perfnone}{HTML}{FFF2CC}
\newcommand{\perfmark}[1]{%
  \ifthenelse{\equal{#1}{good}}{\colorbox{perfgood}{\scriptsize good}}{%
  \ifthenelse{\equal{#1}{poor}}{\colorbox{perfpoor}{\scriptsize poor}}{%
  \ifthenelse{\equal{#1}{noeffect}}{\colorbox{perfnoeffect}{\scriptsize no effect}}{%
  \colorbox{perfnone}{\scriptsize no claim}}}}}

%   \begin{document} ...
%
% Figures and tables are NOT copied: \graphicspath and \reporttable resolve into
% final_report/figures/ and final_report/tables/, so the raw write-up always
% shows the same artefacts as the chapter.  Compile with:  latexmk -pdf <file>

\usepackage[T1]{fontenc}
\usepackage[utf8]{inputenc}
\usepackage{mathptmx}
\usepackage[scaled=0.92]{helvet}
\usepackage{courier}
\usepackage[a4paper,left=30mm,right=20mm,top=25mm,bottom=25mm]{geometry}
\usepackage{microtype}
\usepackage{booktabs}
\usepackage{tabularx}
\usepackage{longtable}
\usepackage{array}
\usepackage{multirow}
\usepackage{amsmath}
\usepackage{amssymb}
\usepackage{ifthen}
\usepackage{siunitx}
\usepackage{graphicx}
\usepackage{placeins}
\usepackage{caption}
\usepackage{enumitem}
\usepackage[table]{xcolor}
\usepackage{rotating}
\PassOptionsToPackage{hyphens}{url}\usepackage{url}
\usepackage[hidelinks]{hyperref}
\usepackage[nameinlink,noabbrev]{cleveref}

\graphicspath{{../../final_report/}{../../final_report/figures/}}
% \reporttable{name} inputs final_report/tables/name.tex
\newcommand{\reporttable}[1]{\input{../../final_report/tables/#1}}

\newcommand{\repo}[1]{\texttt{\detokenize{#1}}}
\sisetup{detect-all,group-separator={,},range-phrase=--,range-units=single}
\DeclareSIUnit{\pixel}{px}
\DeclareSIUnit{\mebi}{Mi}
\captionsetup{font=small,labelfont=bf}
\setlist{nosep}
\setlength{\parskip}{4pt}
\setlength{\LTcapwidth}{\textwidth}

% ---- result-slot status (RENDERED here, unlike in the final report)
\definecolor{stcompleted}{HTML}{A5D6A7}
\definecolor{stpartial}{HTML}{FFF2CC}
\definecolor{stuntested}{HTML}{FFD8A8}
\definecolor{stblocked}{HTML}{FFCDD2}
\definecolor{stfailed}{HTML}{FFCDD2}
\definecolor{stplanned}{HTML}{E0E0E0}
\newcommand{\slotstatus}[2]{%
  \par\noindent\fcolorbox{black!40}{st#1}{\parbox{\dimexpr\textwidth-2\fboxsep-2\fboxrule}{%
    \textbf{Status: #1.}\ #2}}\par\medskip}
\newcommand{\status}[1]{\textbf{#1}}

% ---- queries: RENDERED in red here so a reader sees every open concern in place
\definecolor{flagred}{HTML}{B00000}
\newcommand{\flag}[1]{\textcolor{flagred}{#1}}
\newcommand{\moved}[1]{#1}
\newenvironment{queried}[1][Questionable]%
  {\par\addvspace{0.6ex}\begingroup\color{flagred}\noindent
   \textbf{[#1]}\ \ignorespaces}%
  {\par\endgroup\addvspace{0.6ex}}

% ---- per-block performance marks (section 4.5 diagrams)
\definecolor{perfgood}{HTML}{A5D6A7}
\definecolor{perfpoor}{HTML}{FFCDD2}
\definecolor{perfnoeffect}{HTML}{D9EAF7}
\definecolor{perfnone}{HTML}{FFF2CC}
\newcommand{\perfmark}[1]{%
  \ifthenelse{\equal{#1}{good}}{\colorbox{perfgood}{\scriptsize good}}{%
  \ifthenelse{\equal{#1}{poor}}{\colorbox{perfpoor}{\scriptsize poor}}{%
  \ifthenelse{\equal{#1}{noeffect}}{\colorbox{perfnoeffect}{\scriptsize no effect}}{%
  \colorbox{perfnone}{\scriptsize no claim}}}}}

%   \begin{document} ...
%
% Figures and tables are NOT copied: \graphicspath and \reporttable resolve into
% final_report/figures/ and final_report/tables/, so the raw write-up always
% shows the same artefacts as the chapter.  Compile with:  latexmk -pdf <file>

\usepackage[T1]{fontenc}
\usepackage[utf8]{inputenc}
\usepackage{mathptmx}
\usepackage[scaled=0.92]{helvet}
\usepackage{courier}
\usepackage[a4paper,left=30mm,right=20mm,top=25mm,bottom=25mm]{geometry}
\usepackage{microtype}
\usepackage{booktabs}
\usepackage{tabularx}
\usepackage{longtable}
\usepackage{array}
\usepackage{multirow}
\usepackage{amsmath}
\usepackage{amssymb}
\usepackage{ifthen}
\usepackage{siunitx}
\usepackage{graphicx}
\usepackage{placeins}
\usepackage{caption}
\usepackage{enumitem}
\usepackage[table]{xcolor}
\usepackage{rotating}
\PassOptionsToPackage{hyphens}{url}\usepackage{url}
\usepackage[hidelinks]{hyperref}
\usepackage[nameinlink,noabbrev]{cleveref}

\graphicspath{{../../final_report/}{../../final_report/figures/}}
% \reporttable{name} inputs final_report/tables/name.tex
\newcommand{\reporttable}[1]{\input{../../final_report/tables/#1}}

\newcommand{\repo}[1]{\texttt{\detokenize{#1}}}
\sisetup{detect-all,group-separator={,},range-phrase=--,range-units=single}
\DeclareSIUnit{\pixel}{px}
\DeclareSIUnit{\mebi}{Mi}
\captionsetup{font=small,labelfont=bf}
\setlist{nosep}
\setlength{\parskip}{4pt}
\setlength{\LTcapwidth}{\textwidth}

% ---- result-slot status (RENDERED here, unlike in the final report)
\definecolor{stcompleted}{HTML}{A5D6A7}
\definecolor{stpartial}{HTML}{FFF2CC}
\definecolor{stuntested}{HTML}{FFD8A8}
\definecolor{stblocked}{HTML}{FFCDD2}
\definecolor{stfailed}{HTML}{FFCDD2}
\definecolor{stplanned}{HTML}{E0E0E0}
\newcommand{\slotstatus}[2]{%
  \par\noindent\fcolorbox{black!40}{st#1}{\parbox{\dimexpr\textwidth-2\fboxsep-2\fboxrule}{%
    \textbf{Status: #1.}\ #2}}\par\medskip}
\newcommand{\status}[1]{\textbf{#1}}

% ---- queries: RENDERED in red here so a reader sees every open concern in place
\definecolor{flagred}{HTML}{B00000}
\newcommand{\flag}[1]{\textcolor{flagred}{#1}}
\newcommand{\moved}[1]{#1}
\newenvironment{queried}[1][Questionable]%
  {\par\addvspace{0.6ex}\begingroup\color{flagred}\noindent
   \textbf{[#1]}\ \ignorespaces}%
  {\par\endgroup\addvspace{0.6ex}}

% ---- per-block performance marks (section 4.5 diagrams)
\definecolor{perfgood}{HTML}{A5D6A7}
\definecolor{perfpoor}{HTML}{FFCDD2}
\definecolor{perfnoeffect}{HTML}{D9EAF7}
\definecolor{perfnone}{HTML}{FFF2CC}
\newcommand{\perfmark}[1]{%
  \ifthenelse{\equal{#1}{good}}{\colorbox{perfgood}{\scriptsize good}}{%
  \ifthenelse{\equal{#1}{poor}}{\colorbox{perfpoor}{\scriptsize poor}}{%
  \ifthenelse{\equal{#1}{noeffect}}{\colorbox{perfnoeffect}{\scriptsize no effect}}{%
  \colorbox{perfnone}{\scriptsize no claim}}}}}

\setcounter{chapter}{4}
\renewcommand{\thesection}{\arabic{section}}
\renewcommand{\thefigure}{I.\arabic{figure}}
\renewcommand{\thetable}{I.\arabic{table}}
% every table and figure of this document is generated by script/exp3_insights.py
\newcommand{\insdir}{../../../../output/compare/simulation/exp3-insights}
\newcommand{\ins}[1]{\input{\insdir/#1}}
\newcommand{\patdir}{../../../../output/compare/simulation/exp1-error-pattern}
\newcommand{\pat}[1]{\input{\patdir/#1}}
\newcommand{\realdir}{../../../../output/compare/real}
\newcommand{\insreal}[1]{\input{\realdir/realv6-insights/#1}}
\newcommand{\insrealb}[1]{\input{\realdir/realv7-insights/#1}}

\title{Experiment 1, visual arms: What the Pose Error Rises With\\
       \large Insight analysis for Chapter 4, Section 4.1}
\author{Bolton Athit Davies}
\date{Reporting cutoff: 22 October 2026; real-robot section added 2 November 2026, rebuilt 3 November (realv6) with realv7 added the same day 2026; revised 4 November 2026}

\begin{document}
\maketitle

\noindent\textbf{Purpose of this document.} The raw write-up of Experiment~1
(\path{insights/experiment_1/}) reports what each run scored. The visual arms
analysed here were Experiment~3 until the campaign was restructured on
11 November 2026 and they were merged into Experiment~1; the analysis is
unchanged, only its number. This document asks
\emph{why}: whether the pose error rises when the robot turns, when a movable
object is in the camera's view, or after the tracker loses features, and what
the robot is doing at the moment the error starts to grow. Every question is
tested within each run, so that the run is its own control, and the answer is
stated with the evidence that supports it and the evidence that is missing.
Every number is written by \path{script/exp3_insights.py} into
\path{output/compare/simulation/exp3-insights/}; the tables below are input
from there unchanged.

\section{Questions and method}

Five questions were put to the 114 Experiment~3 runs of both testbeds (the
registry of \path{script/analyse_exp3_timelines.py}: 64 dynamic and 18 static
runs in the extreme world, 26 dynamic and 6 static in the realistic world).

\begin{enumerate}[label=Q\arabic*.]
\item Does the error grow faster while the robot turns than on straights?
\item Does it grow faster while a movable object is in the camera's view?
\item Is the onset preceded by a loss of tracked features?
\item What is the robot doing at the onset: turning or straight, object in view
      or not?
\item Is the error that these questions concern the estimator's error at all,
      or does part of it come from how the estimate is scored?
\end{enumerate}

\paragraph{Within-run contrasts.} The error series is the start-anchored planar
error of the timeline analysis: the estimate is aligned to the ground truth by a
rigid fit over the first \SIrange{10}{40}{\second}, after the initialisation
transient and before the ground truth has turned \SI{60}{\degree}, and the
planar distance to the ground truth is followed from there. For Q1 and Q2 the
samples of a run are split into two states and the growth rate of the error
(metres per second of positive and negative change) is computed in each. States
are: \emph{turning}, ground-truth yaw rate above \SI{10}{\degree\per\second};
\emph{object in view}, for the extreme bags the projected simulator truth of
Experiment~2 (image area of moving objects in the left camera), for the realistic
bags the carriers and carried loads whose poses the bag records on
\path{/world/default/pose/info}, placed within \SI{20}{\metre} and
$\pm\SI{45}{\degree}$ of the ground-truth heading with no shelf footprint on the
line of sight. Rates are taken only while the estimate is still tracking: after
the anchor window, before its first impossible-speed jump and while the error is
below \SI{3}{\metre}. A contrast needs at least \SI{3}{\second} in each state;
across runs the two rates are compared with the sign count and the Wilcoxon
signed-rank test. For Q3 the matched-feature count (ORB-SLAM3 map inliers,
VINS-Fusion tracked features) is followed around the onset as a percentage of
the run's own level before it. For Q4 the state at the onset time is read off.
The onset is the one the raw write-up and the timeline tables use: \SI{1}{\second}
of sustained excess over three times the static median error, with the ground
truth scored as the chapter scores it. The drift contrasts of Q1 and Q2 are
instead computed with the ground truth at the IMU point (Insight~1), because at
the model origin they measure the lever arm, not the estimator.

\paragraph{Lever arm (Q5).} Both estimators report the pose of the IMU: the
ORB-SLAM3 wrapper converts the camera pose with the camera-to-body extrinsic
before writing it, and VINS-Fusion writes its body state. The ground truth is the
model's odometry, which reports the model origin. In the simulated platform the
IMU sits \SI{0.338}{\metre} ahead of that origin
(\path{models/ackermann_robot/model.sdf}). No scoring step in the project
compensates this; Chapter~3 states the requirement only for the real robot's
tracker marker. The effect is tested by moving the ground truth forward along
its own heading by a lever arm $L$ and re-scoring every tracked unfiltered run
for $L$ from \SIrange{0}{0.60}{\metre}.

\FloatBarrier
\section{Insight 1: the error that appears at turns is the scoring lever arm}

Scored as the chapter scores them, the error of the tracked runs grows while the
robot turns and holds on straights; scored at the IMU point, that pattern
disappears and the whole-run ATE falls by half or more.
\Cref{tab:insight-lever,tab:insight-lever-argmin} and \cref{fig:lever} carry
the sweep.

\ins{lever_arm_sweep.tex}
\ins{lever_arm_argmin.tex}

\begin{figure}[htbp]
  \centering
  \includegraphics[width=\textwidth]{\insdir/fig_lever_sweep_light.png}
  \caption{Whole-run planar ATE of every tracked unfiltered run against the
  lever arm assumed between the model origin and the IMU. Solid lines are static
  recordings, dashed the realistic dynamic recordings. The dotted line is the
  value in the robot model.}
  \label{fig:lever}
\end{figure}

The minimum of the ATE lies at the lever arm the robot model specifies. Over the
30 tracked unfiltered runs the ATE-minimising $L^{*}$ has a median of
\SI{0.34}{\metre} and an interquartile range of \SIrange{0.32}{0.40}{\metre}
(\cref{tab:insight-lever-argmin}), for both estimators and both testbeds. The
three runs that do not follow the pattern are the ones whose error is not a
rigid offset: the two ORB-SLAM3 runs on \repo{realistic_002}, whose error is a
\SIrange{0.5}{1.0}{\metre} step taken in the first \SI{15}{\second}
(\cref{fig:traffic}), and the VINS-Fusion run on the same bag, which diverges at
its first turn. With $L=\SI{0.34}{\metre}$ the median static ATE of the extreme
world falls from \SI{0.32}{\metre} to \SI{0.10}{\metre} for VINS-Fusion and from
\SI{0.37}{\metre} to \SI{0.19}{\metre} for ORB-SLAM3; in the realistic world from
\SI{0.37}{\metre} to \SI{0.25}{\metre} and from \SI{0.32}{\metre} to
\SI{0.15}{\metre} (\cref{tab:insight-lever}).

The turning effect goes with it. At $L=0$ ORB-SLAM3's error in the realistic
world grows at a median \SI{8.7}{\milli\metre\per\second} while turning and
shrinks on straights, in eight of eight unfiltered runs; at $L=\SI{0.34}{\metre}$
the turning rate is \SI{0.0}{\milli\metre\per\second}. This is the signature of
a rigid offset between the scored point and the reference point: the offset
vector rotates with the robot, so the apparent error changes only while the
heading changes.

The sweep fits one parameter and could, on its own, absorb some estimator error.
What makes it a finding rather than a fit is that the minimum coincides with an
independently specified geometric constant, in 30 runs, across two estimators
that share no code. The residual ATE at $L^{*}$ is the estimator's.

\begin{queried}[Consequence beyond this experiment]
Every simulation ATE in Chapter~4 that aligns an estimate to the model-origin
ground truth carries this offset: the static floor of Experiment~3 and the
visual estimators of Experiment~1. The wheel and IMU baselines of Experiment~1
are referenced to the same origin and are not affected in the same way, so the
estimator ranking of that experiment has to be re-scored before it is quoted.
The correction is one line: move the ground truth to the IMU point before
alignment. It has since been applied in the scoring scripts
(\path{vio_metrics.load_gt}, 22 October 2026) and every simulation table and
figure of Experiments~1 and~3 has been regenerated; the chapter prose that
quotes the old values is pending.
\end{queried}

\FloatBarrier
\section{Insight 2: after the correction, turning does not raise the error}

With the ground truth at the IMU point, neither estimator's error grows faster
while turning than on straights, in either testbed, and the heading error does
not either. \Cref{tab:insight-contrast-extreme,tab:insight-contrast-realistic}
give the contrasts; \cref{fig:turn} plots each run.

\ins{contrast_extreme_lever0p338.tex}
\ins{contrast_realistic_lever0p338.tex}

\begin{figure}[htbp]
  \centering
  \includegraphics[width=\textwidth]{\insdir/fig_turn_contrast_lever0p338_light.png}
  \caption{Drift rate while turning against drift rate on straights, one marker
  per run, ground truth at the IMU point. Negative rates are drawn at the axis
  floor. A run above the diagonal drifts faster while turning.}
  \label{fig:turn}
\end{figure}

The static recordings of the extreme world are the only group with enough runs
for a test, and the two estimators point opposite ways there: VINS-Fusion is
faster while turning in seven of nine runs and ORB-SLAM3 in one of nine, neither
at $p<0.05$. In the realistic world the sign is mixed in every group. The
heading-error contrast is flat throughout. On these data the answer to Q1 is no.

\FloatBarrier
\section{Insight 3: in the extreme world the failure is immediate and is not
         modulated by how much moving object is in view}

Every dynamic run of the extreme world passes its onset threshold on the first
aisle, almost never while turning, and the moment-to-moment amount of moving
object in the image does not predict when ORB-SLAM3 loses tracking or when
either estimate jumps. \Cref{tab:insight-onset-extreme,tab:insight-mechanism-extreme}
and \cref{fig:epoch} carry the evidence.

\ins{onset_context_extreme.tex}
\ins{mechanism_extreme.tex}

\begin{figure}[p]
  \centering
  \includegraphics[width=\textwidth]{\insdir/fig_onset_epoch_light.png}
  \caption{What precedes the onset, at the chapter's onset definition: median
  over runs (shaded, interquartile range) of the matched-feature count relative
  to the run's own pre-onset level, the true moving-object area in view (extreme
  world) or the number of traffic agents in view (realistic world), the image
  area the filter masked, and the ground-truth yaw rate, against time from
  onset.}
  \label{fig:epoch}
\end{figure}

Moving objects are in view for \SIrange{74}{91}{\percent} of the frames of these
four bags (projected truth of Experiment~2), so a within-bag ``no object''
baseline does not exist; the comparison has to be between more and less moving
area. Split into terciles of moving area within each run, the share of frames
in which ORB-SLAM3 is not tracking OK is not higher in the top tercile than in
the bottom one (\cref{tab:insight-mechanism-extreme}), nor is the jump rate of
either estimate. The matched-feature count falls weakly with moving area
(median Spearman $\rho$ about $-0.1$ for ORB-SLAM3, in eleven of twelve runs).

The onset is on a straight. Of the 64 dynamic runs, the robot was turning at
onset in fifteen (\cref{tab:insight-onset-extreme}); the first heading change of
every route comes after the first aisle, and the error has passed threshold
before it.

The two filtered estimators reach their onset differently. For ORB-SLAM3 with
the class mask the matched features collapse to about one fifth of the run's
own level \SI{2.5}{\second} before the onset (\cref{fig:epoch}, top left): the
mask removes the features the tracker was using, and the error follows.
VINS-Fusion shows no feature loss before its onset, filtered or not; its
estimate does not drift, it breaks, with a jump rate of about 20 per second
once it has failed. What breaks it is not identified by these data; the
hypothesis that wrongly tracked moving points corrupt the visual-inertial
initialisation remains untested.

\FloatBarrier
\section{Insight 4: in the realistic world the traffic in view does not
         measurably change the drift, and the class mask costs features when
         traffic is in view}

With the carriers' true positions from the bag, the drift of the tracked
realistic runs is not consistently faster while a carrier is in the camera's
view, and ORB-SLAM3 with the class mask has fewer matched features during those
moments without a drift penalty.
\Cref{tab:insight-traffic-realistic} and \cref{fig:traffic} carry the
evidence; the per-run values are in
\path{output/compare/simulation/exp3-insights/realistic_runs_lever0p338.tex}.

\ins{traffic_realistic_lever0p338.tex}

\begin{figure}[htbp]
  \centering
  \includegraphics[width=\textwidth]{\insdir/fig_traffic_timeline_lever0p338_light.png}
  \caption{Realistic dynamic recordings, unfiltered runs: anchored planar error
  (ground truth at the IMU point) with the moments a carrier or carried load is
  in the camera's view (violet) and the moments the robot turns (grey). The
  green band is the alignment window.}
  \label{fig:traffic}
\end{figure}

Traffic is in view for a median \SI{37}{\percent} of the tracked time, so both
states are well populated; the groups of runs are not. Two unfiltered
VINS-Fusion runs and four ORB-SLAM3 runs carry the unfiltered answer, and in
them the drift is, if anything, slower while a carrier is in view (median
\SI{-0.2}{\milli\metre\per\second} against \SI{1.5}{\milli\metre\per\second}
for ORB-SLAM3, one run of four faster in view). The filtered ORB-SLAM3 group is
the only one large enough to test: there the matched-feature count is lower
while traffic is in view in eleven of sixteen runs ($p=0.05$) and the drift
rate is higher in seven of ten ($p=0.08$), with no change in the tracking-loss
share. Taken together these suggest, and do not establish, that the class mask
costs features when a carrier is in view and that the cost is paid in drift;
the unfiltered estimators show no such cost.

The one large degradation of the realistic world is not a traffic encounter.
The VINS-Fusion error on \repo{realistic_002}, \SI{1.8}{\metre} whole-run ATE,
is a single \SI{3}{\metre} step taken between \SI{29}{\second} and
\SI{33}{\second}, at the end of the first heading change and with no carrier
in the camera's view (\cref{fig:traffic}, third panel); the carrier that passed
at \SIrange{5}{12}{\second} left no trace. The same bag gives both ORB-SLAM3
runs their \SIrange{0.5}{1.0}{\metre} step at \SI{2.5}{\second} and
\SI{13}{\second}, the moments that same carrier enters and leaves the view.
Whether the step is the estimator's or the ground truth's is not settled by
these data.

\begin{queried}[Coverage of the traffic measurement]
Walking people are gz-sim actors and do not appear on the pose topic, so their
visibility is unmeasured; occlusion is tested against shelf footprints only; and
the carried loads move only while carried. The in-view state is therefore a
lower bound on the moments something moving was in the image.
\end{queried}

\FloatBarrier
\section{Insight 5: real robot (\repo{dataset_realv6_000}): ORB-SLAM3 tracks the
         figure-of-eight to \SI{0.5}{\metre} with the GPU front end, loses a factor
         of two with the CPU front end, and VINS-Fusion has not yet had a valid run}

The first real recording with a usable optical-tracker reference since the
camera remount gives ORB-SLAM3 a provisional planar ATE of \SI{0.46}{} and
\SI{0.52}{\metre} in the two runs with the GPU feature extractor and \SI{0.62}{}
and \SI{1.03}{\metre} in the two with the CPU extractor, which drops two thirds
of the frames at real-time replay; no run loses tracking after its
initialisation resets. VINS-Fusion's run resets 217 times on gaps in its IMU
stream and stops after \SI{19}{\second}. Every number rests on a tracker export
that fails the scorer's validity rules and on one spurious sample whose removal
is documented below. \Cref{tab:realv6-final,tab:realv6-cleaning,tab:realv6-turning,tab:realv6-vins}
and \cref{fig:realv6-paths,fig:realv6-gt,fig:realv6-orb,fig:realv6-vins} carry
the evidence.

\paragraph{The recording.} \repo{dataset_realv6_000} was recorded on 31 October
2026 on the Gensurv floor: \SI{220}{\second}, 6\,601 stereo pairs, the RealSense
IMU at \SI{200}{\hertz}, wheel-speed feedback, and a figure-of-eight route of
about \SI{27}{\metre} at a median \SI{0.15}{\metre\per\second}, no scripted
moving objects.

\paragraph{The runs.} Five runs made on the afternoon of 3 November 2026 are
reported, the only ones whose run directories still exist: ORB-SLAM3
stereo-inertial with the GPU or the CPU feature extractor and with loop closing
on or off, and VINS-Fusion once; all at replay rate 1.0, all without a detector,
so unfiltered. ORB-SLAM3 used \repo{config/wil/stereo_imu.yaml} with the
repaired library (fork patches 11--14b, node and library built alike), the
3 November stereo extrinsic and \texttt{IMU.NoiseGyro} 0.1 with
\texttt{IMU.NoiseAcc} 1.0, ten times the values that had been copied from the
VINS file; VINS-Fusion used the same file's gravity norm of \num{9.2642} from an
August recording. Earlier runs on this bag (2 November, and the morning of
3 November) were overwritten by these re-runs; their scorings remain under
\path{output/compare/real/realv6*} as a record of the search described in
\cref{sec:realv6-stability} but cannot be reproduced and carry no number here.
No filtered run exists, so the filter hypothesis remains untested on the real
robot.

\paragraph{Reporting rules.} Scoring is \path{script/eval_real_gt.py}: rigid
SE(2) alignment with scale fixed, lever arm from the tracker marker to the IMU
unmeasured and set to zero, every number \emph{provisional}. The bag was not on
the machine when these runs were scored, so the time offset is pinned to the
value the wheel-speed cross-correlation found on 2 November
(\SI{-4.14}{\second} on the cleaned export, \SI{+4.38}{\second} on the raw one)
and the bag's first-message time and duration are taken from that day's
manifest; the new manifests record the substitution. The scorer's validity
rules were fixed before the data were read; the tracker export fails two of them
(gaps above \SI{1}{\second}, a step above \SI{0.5}{\metre}), so the tool labels
every run \emph{failed (ground truth invalid)} and the numbers are reported as
diagnostics, not as a result-slot headline. Outputs:
\path{output/compare/real/realv6_1103pm_*/} and
\path{output/compare/real/realv6-insights/}.

\insreal{gt_cleaning.tex}

\begin{figure}[htbp]
  \centering
  \includegraphics[width=\textwidth]{\realdir/realv6-insights/fig_gt_cleaning_light.png}
  \caption{The tracker export of \repo{dataset_realv6_000}: one sample at
  \SI{184.9}{\second} lies \SI{5.2}{\metre} off the path and back (left), which
  puts an \SI{8.7}{\metre\per\second} spike into the speed profile the time
  alignment correlates against (right).}
  \label{fig:realv6-gt}
\end{figure}

One tracker sample decided the headline. The export runs at \SI{1.7}{\hertz},
not the \SI{7}{\hertz} of the earlier recordings, with 66 gaps longer than
\SI{1}{\second}, and it contains a single sample \SI{5.2}{\metre} off the path.
That sample adds \SI{10}{\metre} to the true path length (37.3 against
\SI{27.1}{\metre}) and dominates the speed cross-correlation; scored against it
the GPU run reads \SI{1.17}{\metre} with a path ratio of 0.65, and with that one
sample removed, and nothing else, \SI{0.52}{\metre} with a path ratio of 0.85
(\cref{tab:realv6-cleaning}). The sparse rate and the gaps remain, so the cleaned
truth still fails the gap rule; the cleaning is recorded here as a one-sample
edit with its effect shown, not applied silently.

\begin{figure}[p]
  \centering
  \includegraphics[width=\textwidth]{\realdir/realv6-insights/fig_paths_light.png}
  \caption{Every surviving run against the cleaned tracker truth after rigid
  SE(2) alignment, one panel per run: ORB-SLAM3 with GPU or CPU extractor and
  loop closing on or off, and VINS-Fusion, whose \SI{18}{\second} before its
  resets took over are all it tracked.}
  \label{fig:realv6-paths}
\end{figure}

\insreal{final_table.tex}

\begin{figure}[htbp]
  \centering
  \includegraphics[width=\textwidth]{\realdir/realv6-insights/fig_orb_timeline_light.png}
  \caption{ORB-SLAM3 planar error against the cleaned tracker truth over the
  run, one panel per configuration. Grey bands: the aligned estimate's heading
  changes faster than \SI{10}{\degree\per\second}. Dotted red: the frames at
  which tracking left the OK state (the \repo{local_mapping_bad_imu}
  initialisation resets, all inside the first \SI{13}{\second}).}
  \label{fig:realv6-orb}
\end{figure}

\insreal{turning.tex}

ORB-SLAM3 tracks the whole route in all four configurations, and the front end
is what separates them. The two GPU runs score \SI{0.52}{} (loop closing on) and
\SI{0.46}{\metre} (off, with zero resets) after dropping 2\,420 and 3\,201 of
about 6\,480 frames; the two CPU runs drop 4\,363 and 4\,204 frames, two thirds
of the stream, and score \SI{0.62}{} and \SI{1.03}{\metre}, the first with a
\SI{2}{\metre} vertical excursion the scorer flags as implausible and the second
with three resets and a half-second of lost tracking
(\cref{tab:realv6-final}, \cref{fig:realv6-paths}). Loop closing changes nothing
that the front end does not change more. After the initialisation resets no run
leaves the OK state; the error is a slow wave between \SI{0.05}{} and
\SI{1}{\metre}, not an accumulating drift (\cref{fig:realv6-orb}), and it does
not rise while turning in any of the four (\cref{tab:realv6-turning}), the same
answer the simulation gave once the lever arm was removed.

What the \SI{0.5}{\metre} is made of is not settled. The estimates' loops are a
uniform \SIrange{12}{15}{\percent} smaller than the tracker's (path ratios
0.85--0.88 in the three plausible runs; \cref{fig:realv6-paths}), which a
stereo-baseline error, a tracker scale error or the unmeasured lever arm could
each produce.

\subsection{How the setting was reached}
\label{sec:realv6-stability}

The configuration above was not the first tried. The runs of 2 November and of
the morning of 3 November, whose directories these re-runs overwrote, showed
that the stereo-inertial estimator with the IMU noise densities copied from the
VINS file was unstable on this recording: a minority of runs stayed below
\SI{1}{\metre} and the rest ran away in the velocity state at a different
place each time, while the robot crept at \SI{0.12}{\metre\per\second}. Six
out-of-bounds reads in the stereo-fisheye code path were found and fixed, an
AddressSanitizer run over the whole recording came out clean, and keyframe
culling, loop closing, the GPU front end, the replay rate, frame drops, the
stereo extrinsic and the initialisation were each varied without removing the
runaway; raising the IMU noise densities tenfold removed it in every run since.
The reading is that the inertial optimisation was trusting an IMU whose
camera-relative axis signs, lever arm and time offset are unverified; with the
weaker weighting the images carry the solution and the IMU acts as a prior,
which is the role the evidence supports. The search is recorded in the ORB-SLAM3
engineering record (sections 7.10 to 7.13 and the status table); no number from
those overwritten runs is quoted here.

\insreal{vins_table.tex}

\begin{figure}[htbp]
  \centering
  \includegraphics[width=\textwidth]{\realdir/realv6-insights/fig_vins_light.png}
  \caption{VINS-Fusion on \repo{dataset_realv6_000}: distance of the estimate
  from its origin with every \repo{imu_gap} reset marked (top), and the IMU
  messages the node received against the samples the estimator used (bottom).}
  \label{fig:realv6-vins}
\end{figure}

VINS-Fusion has no valid run on this recording. The node received 3\,187 of the
bag's 6\,601 frames but used only 6\,498 of the 43\,494 IMU messages it
received; it raised 217 \repo{imu_gap} resets and wrote \SI{19}{\second} of
poses before the resets took over (\cref{tab:realv6-vins}, \cref{fig:realv6-vins}).
The delivery of the IMU stream into VINS-Fusion is therefore the defect in front
of every other; behind it stands a gravity norm of \num{9.2642} in the
configuration against the \SI{9.878}{\metre\per\second\squared} this bag's IMU
reads at rest, which the overwritten 2 November runs had shown to drive a
velocity runaway once the visual constraints thin. The \SI{18}{\second} it did
track scored \SI{0.18}{\metre} against 41 tracker samples, which is too short to
mean anything and is not a result.

\subsection{realv7: three recordings with a valid tracker truth, and an IMU
            initialisation that the map reset threw away}
\label{sec:realv7-tilt}

The three \repo{dataset_realv7} recordings of 3 November 2026 are the first
whose tracker export passes every validity rule of the scorer:
\SIrange{5}{6}{\hertz}, no gap above \SI{1}{\second}, no step. Their routes
are \SI{21.6}{}, \SI{31.7}{} and \SI{23.4}{\metre} over \SI{97}{}, \SI{134}{}
and \SI{104}{\second}: a rounded rectangle, a figure-of-eight and a rounded
rectangle.

On them the ORB-SLAM3 setting that tracked realv6 writes an estimate that lies
in a tilted plane in nine of twelve runs (\cref{tab:realv7-tilt},
\cref{fig:realv7-tilt}); scored as written, those nine runs are not scorable.
The tilt was read first as a wrong gravity direction and then as a frame fixed
by the node; the per-frame diagnostics read on 4 November give the mechanism,
and a repair in the scorer that takes nothing from the truth makes all twelve
runs scorable, with the time offsets from the wheel speed now that the bags are
on the dataset drive (\cref{tab:realv7-results},
\cref{fig:realv7-paths-repaired}).

\insrealb{tilt_table.tex}

\begin{figure}[p]
  \centering
  \includegraphics[width=\textwidth]{\realdir/realv7-insights/fig_tilt_light.png}
  \caption{The twelve \repo{realv7} ORB-SLAM3 estimates as written in the live
  log, seen from the side (estimate $x$ against $z$, centred). A floor robot's
  path should be a flat line; on \repo{realv7_000} every run is the route drawn
  in a vertical plane. The repaired estimates are in
  \cref{fig:realv7-paths-repaired}.}
  \label{fig:realv7-tilt}
\end{figure}

The tilt is the inertial initialisation being thrown away, not the gravity
estimate and not the node. A plane fitted to each estimate's positions lies
within \SI{10}{\degree} of horizontal in three runs and at
\SIrange{88}{90}{\degree} in five, the rest between (\cref{tab:realv7-tilt});
on \repo{realv7_000} every configuration draws the route in a vertical plane
(\cref{fig:realv7-tilt}, top row), and the plane normal is in every run exactly
one body axis. The node writes ORB-SLAM3's world frame as it is, its first-pose
alignment option being off.

The per-frame diagnostics tell one story in nine runs of twelve. The IMU
initialised \SIrange{3}{4}{\second} after the first pose, and ORB-SLAM3 rotated
its map so that gravity points down $z$; the \SI{0.38}{\metre} step in the log
is that rotation applied to a robot \SI{0.38}{\metre} from the origin. Within
two seconds the upstream rule that resets a map whose consecutive keyframes lie
closer than \SI{2}{\centi\metre} in the ten seconds after the initialisation
fired, because the robot was still standing or crawling, and the map was
discarded.

In the new map the initialisation was starved. That map starts in the camera
frame of the reset instant, whose vertical is the camera's $y$ axis. Before the
IMU is initialised a keyframe is inserted every \SI{0.25}{\second}, creating
map points for one takes about \SI{250}{\milli\second} on this rig, and
upstream attempts the initialisation only when its keyframe queue is empty,
which it was for \SIrange{3}{9}{\percent} of the keyframes. The IMU
re-initialised after \SIrange{37}{92}{\second} in seven runs and never in two
(\repo{000} without loop closing and \repo{001} with it, both GPU front end);
only \repo{001} with the CPU front end and loop closing re-initialised early, at
\SI{6.7}{\second}, and it is the run that was level.

Until the re-initialisation every pose is a stereo-only estimate in a map whose
vertical is the camera's $y$ axis, which is the body-axis normal. At the
re-initialisation the map is rotated about its origin under the already-written
poses, which is the \SIrange{1.9}{6.8}{\metre} step of the late-initialising
runs. The IMU noise weighting that stopped the realv6 runaway and the
IMU-to-camera rotation, whose axis signs remain unverified, play no part. The
four surviving realv6 runs of \cref{tab:realv6-final} show the same reset at
\SIrange{3}{5}{\second} but re-initialised after \SIrange{9}{34}{\second},
which is why they were level.

Repaired, the twelve runs score (\cref{tab:realv7-results},
\cref{fig:realv7-paths-repaired}). The scorer's two new switches take nothing
from the truth. One rotates the poses written before each initialisation by the
rotation measured at the initialisation frame, which brings the \SI{0.38}{} and
\SI{6.8}{\metre} steps to \SI{0.00}{} and \SI{0.03}{\metre}; the other levels
each map's own fitted plane, the side chosen by the body's up axis, which agrees
with the fitted normal within \SI{2}{\degree} wherever the plane is determined.

On \repo{000} the four configurations give \SIrange{0.25}{0.27}{\metre} planar
ATE over \SI{89}{\second}, an RPE at \SI{1}{\second} of \SI{0.05}{\metre} and
path ratios of 1.05 to 1.09, the best real-robot numbers of the project; with
the IMU initialised for at most a quarter of the frames they are stereo-only
figures in all but name. On \repo{001} the four give
\SIrange{0.44}{0.82}{\metre} over \SI{126}{\second}; the run with the IMU
initialised for \SI{94}{\percent} of the frames scores \SI{0.44}{\metre}, no
better than the run with \SI{3}{\percent}. The \SI{0.56}{\metre} run (CPU front
end, loop closing off) keeps a vertical scatter of \SI{1.02}{\metre} after
levelling, two centimetres over the scorer's vertical-range rule, which labels
it failed.

On \repo{002} one run tracks (\SI{0.47}{\metre}, CPU front end, loop closing).
Two lose tracking inside the route (resets at \SI{49}{} and \SI{59}{\second},
\SI{2.1}{\metre} ATE) and the fourth drifts without a reset (\SI{1.67}{\metre},
RPE \SI{0.70}{\metre}, vertical range \SI{4.3}{\metre}, failed by the same
rule). The init-time resets make every other run \emph{partial} under the
scorer's rules, as the table says.

The remedy in the estimator is fork patch 15: two switches in the local mapper
that attempt the initialisation while keyframes are queued and skip the
insufficient-motion reset for stereo, whose scale is fixed. They are on in the
real-rig configuration and untested on a recording; every number above was
obtained without them.

\insrealb{results_table.tex}

\begin{figure}[p]
  \centering
  \includegraphics[width=\textwidth]{\realdir/realv7-insights/fig_paths_lvl_light.png}
  \caption{The twelve \repo{realv7} ORB-SLAM3 runs after the scorer's frame
  repair, against the tracker truth after rigid SE(2) alignment; each panel's
  title gives the planar ATE and the share of tracked frames with the IMU
  initialised.}
  \label{fig:realv7-paths-repaired}
\end{figure}

\insrealb{vins_table.tex}

VINS-Fusion fails on all three in the manner of the first realv6 runs: the IMU
stream now arrives without gaps (zero resets), \SIrange{40}{45}{\percent} of the
frames reach the node, and the estimate diverges to kilometres on the stale
gravity norm.

\begin{queried}[What this slot still needs before it is a result]
For realv6, a tracker export at its nominal \SI{7}{\hertz} without gaps, or an
explicit rule for scoring on a sparse reference; for realv7, one run with the
fork's two initialisation switches on, to show from the per-frame log that the
IMU stays initialised from the first seconds without the scorer's repair; for
both recording sets, the marker-to-IMU lever arm measured and the IMU-to-camera axis
signs verified, which remain open although they no longer explain the tilt; a
VINS-Fusion run whose node receives the IMU stream without gaps and the full
frame stream, with \repo{g_norm} set to what this IMU measures; repeats of the
ORB-SLAM3 configurations kept under new names rather than overwritten; and at
least one filtered run per estimator, without which Experiment~3's hypothesis
has no real-robot test.
\end{queried}

\FloatBarrier
\section{Insight 6: the two estimators err in perpendicular directions, and the
scalar ATE hides it}

VINS-Fusion drifts sideways and ORB-SLAM3 drifts forwards. Decomposing each
run's error into the ground-truth heading frame
(\pat{translation.tex}, \cref{fig:exp1-translation}), VINS-Fusion's cross-track
share of the planar error is \numrange{0.67}{0.73} on three of the four
recordings, while every ORB-SLAM3 run that held tracking lies between
\num{0.11} and \num{0.25}. The medians over the twenty runs are
\SI{0.12}{\metre} along against \SI{0.14}{\metre} across for VINS-Fusion, and
\SI{0.28}{} against \SI{0.12}{\metre} for ORB-SLAM3.

The reported ATE cannot show this. On \repo{realistic_000} ORB-SLAM3 with the
GPU front end scores \SI{0.293}{\metre} and VINS-Fusion \SI{0.290}{\metre} ---
indistinguishable, and perpendicular. For a \num{0.83} by \SI{0.50}{\metre}
robot in a \SI{3.0}{\metre} aisle the two are not equally safe: cross-track
error closes the gap to the racking and along-track error does not.

The decomposition is computed on the same aligned series the reported ATE comes
from, so the parts account for the whole; the proprioceptive family has carried
the same split, anchored rather than aligned, since the raw write-up was
first written, and the two conventions are never mixed in one figure.

\begin{figure}[htbp]
  \centering
  \includegraphics[width=\textwidth]{\patdir/fig_translation_light.png}
  \caption[What the visual estimators' translation error is made of.]{What the
  translation error is made of, on the aligned series the reported ATE uses.
  (a) the error split per run into along-track, cross-track and vertical;
  (b) every sample in the body frame, clipped to the 99th percentile because
  the five diverged runs reach metres; (c) the cross-track share of the planar
  error, where \num{0.5} is isotropic; (d) error against distance travelled,
  on a symmetric-log axis above \SI{0.5}{\metre} so that the eighteen runs that
  held tracking stay visible.}
  \label{fig:exp1-translation}
\end{figure}

\section{Insight 7: when ORB-SLAM3 fails, it fails sideways}

The false closure throws the estimate across the aisle, not along it. The two
loop-closing runs on \repo{realistic_001} are the only ORB-SLAM3 runs whose
cross-track share exceeds \num{0.5} (\num{0.62} and \num{0.69}), and their
cross-track RMS is \SI{4.76}{} and \SI{3.98}{\metre} against along-track
\SI{3.69}{} and \SI{2.64}{\metre}.

That is the failure this experiment already identifies, seen from a new angle:
the \SI{14}{\metre} correction both front ends accept at \SI{73}{\second} does
not merely enlarge the error, it rotates it into the direction that collides.
Of the five diverged runs only these two are cross-track dominated, so the
effect belongs to the false closure rather than to divergence in general.

\section{Insight 8: the error is scatter, not a constant offset}

In simulation the error has no appreciable bias: the mean of each signed
component is \numrange{0.17}{0.36} of its RMS across every arm, so the estimate
scatters about the truth rather than sitting displaced from it.

This is the opposite of what the navigation slot finds on the same maps, where
the error is dominated by a constant offset that the pose inherits from a
locally displaced prior. The two are consistent once the references are
separated: this experiment scores against ground truth directly, so only the
estimator's own jitter is visible, while navigation scores a pose pulled onto a
map that is itself wrong. The contrast is worth carrying into the
cross-experiment discussion, because it says the offset belongs to the map and
the scatter to the estimator.

\section{Insight 9: rotation --- turning does not raise it, and the RMS is one
event on all three axes}

Turning does not raise the rotation error. Over the twenty runs the median
rotation error while turning exceeds the median while driving straight by
\SI{+0.01}{\degree}, and only eleven of twenty runs are positive at all
(\pat{rotation.tex}, \cref{fig:exp1-rotation}); the threshold is a ground-truth
turn rate above \SI{10}{\degree\per\second}, which covers \SI{27}{\percent} of
the samples. This settles for rotation what Insight~2 settled for position.

ORB-SLAM3's \SIrange{14}{19}{\degree} rotational RMS is confirmed as a single
event rather than drift, and the per-axis split says what kind. Yaw, pitch and
roll RMS are all \SIrange{14}{19}{\degree} while their medians are
\SIrange{0.88}{1.51}{\degree}: an error that moves all three axes together and
then persists in none of them is a frame re-alignment, which is what the
inertial initialisation performs.

VINS-Fusion carries a smaller signature of the same kind, and nothing in the
report has yet noticed it. Its pitch and roll RMS are \SI{1.05}{} and
\SI{1.33}{\degree} against a yaw RMS of \SI{0.24}{\degree} --- four to five
times larger out of plane than in it, on a robot that never leaves the floor.
\flag{Whether this is a residual gravity-alignment error or an artefact of the
rigid $SE(3)$ fit absorbing a tilt has not been separated; it is reported as an
observation, not a defect.}

\begin{figure}[htbp]
  \centering
  \includegraphics[width=\textwidth]{\patdir/fig_rotation_light.png}
  \caption[The orientation error by axis and against turn rate.]{The
  orientation error. (a) yaw, pitch and roll RMS per run on a logarithmic
  axis; (b) median rotation error binned by the ground truth's own turn rate;
  (c) the same two medians paired within each run; (d) one run's rotation error
  against time with the turning segments shaded.}
  \label{fig:exp1-rotation}
\end{figure}

\section{What these data cannot answer}

\begin{itemize}
\item \textbf{Whether the filter helps.} No admissible filtered-versus-unfiltered
      pair exists in either simulated testbed, and no filtered run exists on the
      real robot; nothing here changes that.
\item \textbf{Why VINS-Fusion fails in the extreme world.} The failure is
      immediate and not preceded by feature loss or by a rise in moving-object
      area; the mechanism is not identified.
\item \textbf{The traffic effect with confidence.} Two to four unfiltered runs
      per estimator in the realistic world, with the people unmeasured, bound
      the effect rather than estimate it.
\item \textbf{The feature-collapse signal, robustly.} The onset threshold is
      three times the static median error, and that median carries the lever
      arm. Recomputing the thresholds at the IMU point lowers them, moves the
      onsets earlier, and weakens the ORB-SLAM3+YOLO feature collapse of
      \cref{fig:epoch} to a dip to about half in four runs
      (\path{fig_onset_epoch_lever0p338_light.png}). The collapse is real at
      the chapter's definition and threshold-dependent in size.
\end{itemize}

The direction and rotation decompositions of Insights~6 to~9 inherit the slot's
own limit: one run per cell, and ORB-SLAM3's run-to-run spread over identical
settings reaches \SI{0.277}{\metre}, which is wider than most differences
between settings. They are therefore read as patterns pooled over the twenty
runs, never as a ranking of one run against another, and no figure in them
orders runs. The split itself is robust to that limit in a way a difference of
means is not: a cross-track share near \num{0.7} for one estimator and near
\num{0.2} for the other holds on every recording and in both front ends, so it
survives a noise floor that would swallow a \SI{0.1}{\metre} difference in ATE.

\section{Summary}

\begin{table}[htbp]
\centering\small
\caption{Questions, answers and the strength of the evidence.}
\label{tab:insight-summary}
\begin{tabularx}{\textwidth}{lXXl}
\toprule
 & Answer & Evidence & Confidence \\
\midrule
Q5 lever arm & Yes: \SI{0.34}{\metre} of forward offset is in every simulation ATE & minimum of the sweep at the SDF value in 30 runs, two estimators, two testbeds & high \\
Q1 turning & No, once the lever arm is removed & contrasts of \cref{tab:insight-contrast-extreme,tab:insight-contrast-realistic}; heading error flat & medium (small groups) \\
Q2 object in view, extreme & Not measurable within a bag; moving area does not modulate tracking loss or jumps & \cref{tab:insight-mechanism-extreme} & medium \\
Q2 traffic in view, realistic & No effect on the unfiltered drift; the class mask loses features and gains drift while traffic is in view (suggestive) & \cref{tab:insight-traffic-realistic}, \cref{fig:traffic} & low (n) \\
Q3 features before onset & ORB-SLAM3+YOLO yes (collapse to one fifth); VINS-Fusion no & \cref{fig:epoch} & medium \\
Q4 state at onset & On the first straight, not turning; objects in view & \cref{tab:insight-onset-extreme} & high \\
Real robot, ORB & Tracks the whole \SI{218}{\second} route in all four configurations; provisional ATE \SI{0.46}{} and \SI{0.52}{\metre} with the GPU extractor, \SI{0.62}{} and \SI{1.03}{\metre} with the CPU extractor (two thirds of frames dropped); no turning effect; \SIrange{12}{15}{\percent} scale shortfall unexplained & \cref{tab:realv6-final,fig:realv6-paths} & medium (provisional, sparse truth, one run each) \\
Real robot, realv7 & Valid truth; the tilt in 9 of 12 runs is the IMU initialisation thrown away by the upstream low-motion reset and starved in the new map; repaired without the truth, ORB scores 0.25--0.27 m (000), 0.44--0.82 m (001), 0.47 m (002, one run), with two losses and one drift on 002 & \cref{tab:realv7-results,fig:realv7-paths-repaired} & high (mechanism), medium (ATE) \\
Real robot, VINS & No valid run: 217 \repo{imu_gap} resets, \SI{19}{\second} of poses; gravity norm still stale & \cref{tab:realv6-vins}, \cref{fig:realv6-vins} & high \\
Real robot, truth & One tracker sample set the raw ATE to \SI{1.18}{\metre}; truth at \SI{1.7}{\hertz} fails the gap rule & \cref{tab:realv6-cleaning} & high \\
\bottomrule
\end{tabularx}
\end{table}

\paragraph{Consequences for Chapter 4.} Re-score every simulation ATE with the
ground truth moved to the IMU point before it is quoted, in Experiments~1 and~3.
Report the extreme world as an immediate-failure testbed in which the filter
cannot be evaluated, not as a drift testbed. Report the realistic world's
traffic effect as bounded, not measured, and the class mask's feature cost as
the one mechanism these data show. For the real-robot slot, quote ORB-SLAM3's
\SIrange{0.46}{0.52}{\metre} (GPU front end) as provisional with the sparse-truth and lever-arm caveats
attached, report the VINS-Fusion failure as a delivery defect (IMU-stream gaps)
in front of a configuration defect (stale gravity norm) rather than an
estimator result, and keep the slot \emph{partial}
until a filtered run and a full-rate truth exist; report realv7 as the
recording set with valid truth, ORB-SLAM3 at \SIrange{0.25}{0.27}{\metre} on
\repo{000}, \SIrange{0.44}{0.82}{\metre} on \repo{001} and \SI{0.47}{\metre}
on the one \repo{002} run without a loss, after the scorer's frame repair and
with each run's initialisation history beside the number; the IMU stayed
uninitialised for most of nine runs, so these are largely stereo-only figures
until the fork's initialisation switches are exercised.

\end{document}
